\FifthArticle{2}
Das gesamte Gebiet der komplexen Größen, in dem die Werte von $x$ gesucht werden sollen, die der Gleichung $X=0$ genügen, ist ein unendliches zweidimensionales Gebiet. Setzt man nämlich einen solchen Wert in der Form $x=t+iu$ an, wobei $i$ stets die imaginäre Einheit $\surd{-1}$ bezeichnet, so sind für $t$ und $u$ alle reellen Werte von $-\infty$ bis $+\infty$ zulässig. Wir müssen nun zunächst aus diesem unendlichen Gebiet ein begrenztes endliches absondern, außerhalb dessen sicher keine Wurzel der gegebenen Gleichung $X=0$ liegen kann. Dies kann auf mehr als eine Weise geschehen; die folgende scheint unserem Zweck am besten zu entsprechen.

Anstelle der Form $t+iu$ verwende man
\[
x=r(\cos\rho+i\sin\rho)
\]
Um das gesamte unendliche Gebiet der komplexen Größen zu erfassen, muss dann $r$ alle positiven Werte von $0$ bis $+\infty$ und $\rho$ alle Werte von $0$ bis $360^0$ durchlaufen, oder, was
\SourceLeaf{083}
dasselbe ist, von einem beliebigen Anfangswert bis zu einem um $360^0$ größeren Endwert.

Um für $r$ eine Schranke zu erhalten, oberhalb derer kein Wert mehr einer Wurzel der Gleichung $X=0$ entsprechen kann, bringe ich zunächst die Koeffizienten der einzelnen Glieder von $X$ in eine ähnliche Form wie $x$, nämlich
\[
\begin{array}{l}
A=a(\cos\alpha+i\sin\alpha)\\
B=b(\cos\beta+i\sin\beta)\\
C=c(\cos\gamma+i\sin\gamma)\quad\mbox{usw.}
\end{array}
\]
Dabei sollen $a,b,c$ bestimmte positive Größen bezeichnen, abgesehen davon, dass auch die eine oder andere von ihnen $=0$ sein kann. Ich betrachte sodann die Gleichung
\[
r^n-\surd2\,.(ar^{n-1}+br^{n-2}+cr^{n-3}+\mbox{usw.})=0
\]
die, wie man leicht sieht, eine positive Wurzel hat, und zwar nach dem Satz von \textsc{Harriot} nur eine solche. Es sei $R$ diese Wurzel. Dann ist unmittelbar klar, dass für jeden positiven Wert von $r$, der größer als $R$ ist, der Wert von
$r^n-\surd2\,.(ar^{n-1}+br^{n-2}+cr^{n-3}+\mbox{usw.})$
positiv ist und dasselbe auch für die Funktion
\[
nr^n-\surd2\,.((n-1)ar^{n-1}+(n-2)br^{n-2}+(n-3)cr^{n-3}+\mbox{usw.})
\]
gilt, da diese das $n$-Fache der ersten Funktion um
\[
\surd2\,.(ar^{n-1}+2br^{n-2}+3cr^{n-3}+\mbox{usw.})
\]
also um eine positive Differenz übertrifft.

\Needspace{13\baselineskip}\FifthArticle{3}
Ich behaupte nun, dass die Größe $R$ als eine solche Schranke für die Werte von $r$ dienen kann, wie sie im vorhergehenden Artikel verlangt wurde. Der Satz ist folgendermaßen zu beweisen.

\Needspace{10\baselineskip}Ich setze allgemein $X=T+iU$, wobei $T$ und $U$ selbstverständlich reelle Größen bezeichnen. Dabei gilt
\[
\begin{array}{r@{}l}
T={}&r^n\cos n\rho+ar^{n-1}\cos((n-1)\rho+\alpha)+br^{n-2}\cos((n-2)\rho+\beta)\\
&{}+cr^{n-3}\cos((n-3)\rho+\gamma)+\mbox{usw.}
\end{array}
\]
\SourceLeaf{084}
\[
\begin{array}{r@{}l}
U={}&r^n\sin n\rho+ar^{n-1}\sin((n-1)\rho+\alpha)+br^{n-2}\sin((n-2)\rho+\beta)\\
&{}+cr^{n-3}\sin((n-3)\rho+\gamma)+\mbox{usw.}
\end{array}
\]
\Needspace{15\baselineskip}Man sieht leicht, dass, wenn für $r$ irgendein positiver Wert größer als $R$ gewählt wird, $T$ notwendig dasselbe Vorzeichen wie $\cos n\rho$ hat, sobald der Betrag dieses Kosinus nicht kleiner als $\surd\frac12$ ist. Dazu braucht man nur $T$ in die folgende Form zu bringen:
\[
\begin{array}{r@{}l}
\pm T={}&\surd\frac12\,.r^n-ar^{n-1}-br^{n-2}-cr^{n-3}-\mbox{usw.}\\
&{}+[\pm\cos n\rho-\surd\frac12]r^n\\
&{}+[1\pm\cos((n-1)\rho+\alpha)]ar^{n-1}\\
&{}+[1\pm\cos((n-2)\rho+\beta)]br^{n-2}\\
&{}+[1\pm\cos((n-3)\rho+\gamma)]cr^{n-3}\\
&{}+\mbox{usw.}
\end{array}
\]
Hier gelten die oberen Vorzeichen für einen positiven und die unteren für einen negativen Wert von $\cos n\rho$. Der erste Teil des Ausdrucks auf der rechten Seite ist nach dem Satz des vorhergehenden Artikels positiv, während die übrigen Teile zumindest nicht negativ werden können. Auf ganz ähnliche Weise folgt, indem man in der obigen Formel nur $U$ statt $T$ und durchgehend Sinus statt Kosinus schreibt, dass unter derselben Voraussetzung über $r$ stets $U$ dasselbe Vorzeichen wie $\sin n\rho$ hat, sobald der Betrag dieses Sinus nicht kleiner als $\surd\frac12$ ist. Daher hat in allen Fällen mindestens eine der beiden Größen $T,U$ ein im Voraus bestimmtes positives oder negatives Vorzeichen, und folglich kann für keinen Wert von $\rho$ die Gleichung $X=0$ erfüllt sein. W.~Z.~B.~W.

\Needspace{16\baselineskip}\FifthArticle{4}
Um das Verhalten von $T$ und $U$ hinsichtlich ihrer Vorzeichen und Vorzeichenwechsel bei einem festen, $R$ überschreitenden Wert von $r$ noch deutlicher zu machen, lasse man $\rho$ alle Werte zwischen zwei um $360^0$ verschiedenen Grenzen durchlaufen. Dafür sollen jedoch nicht $0$ und $360^0$ gewählt werden, sondern, wenn zur Abkürzung
\[
\frac{45^0}{n}=\omega
\]
\SourceLeaf{085}
gesetzt wird, $-\omega$ und $(8n-1)\omega$. Den gesamten Zwischenraum teile ich in $4n$ gleiche Teile, sodass der erste von $-\omega$ bis $\omega$, der zweite von $\omega$ bis $3\omega$, der dritte von $3\omega$ bis $5\omega$ usw. reicht. Zunächst sind auch die Werte der Ableitungen $\frac{\mathrm{d}T}{\mathrm{d}\rho}$, $\frac{\mathrm{d}U}{\mathrm{d}\rho}$ zu betrachten, für die gilt
\[
\begin{array}{rl}
\frac{\mathrm{d}T}{\mathrm{d}\rho}
&=-nr^n\sin n\rho-(n-1)ar^{n-1}\sin\bigl((n-1)\rho+\alpha\bigr)\\
&\qquad -(n-2)br^{n-2}\sin\bigl((n-2)\rho+\beta)\bigr)\\
&\qquad -(n-3)cr^{n-3}\sin\bigl((n-3)\rho+\gamma\bigr)-\mathrm{usw.}\\
\frac{\mathrm{d}U}{\mathrm{d}\rho}
&=nr^n\cos n\rho+(n-1)ar^{n-1}\cos\bigl((n-1)\rho+\alpha\bigr)\\
&\qquad +(n-2)br^{n-2}\cos\bigl((n-2)\rho+\beta)\bigr)\\
&\qquad +(n-3)cr^{n-3}\cos\bigl((n-3)\rho+\gamma\bigr)+\mathrm{usw.}
\end{array}
\]
Daraus erkennt man leicht durch ähnliche Schlüsse wie im vorhergehenden Artikel und unter Verwendung des Satzes am Ende von Artikel 2, dass $\frac{\mathrm{d}T}{\mathrm{d}\rho}$ stets das entgegengesetzte Vorzeichen von $\sin n\rho$ hat, sobald der Betrag dieses Sinus nicht kleiner als $\surd\frac{1}{2}$ ist, während $\frac{\mathrm{d}U}{\mathrm{d}\rho}$ stets dasselbe Vorzeichen wie $\cos n\rho$ hat, sobald der Betrag dieses Kosinus nicht kleiner als $\surd\frac{1}{2}$ ist. Daraus ergeben sich die folgenden Schlüsse.

\Needspace{14\baselineskip}Im ersten Intervall, also von $\rho=-\omega$ bis $\rho=+\omega$, ist $T$ stets positiv. Dagegen ist $U$ beim Anfangswert negativ und beim Endwert positiv, also dazwischen sicher einmal $=0$, und zwar \emph{nur} einmal, weil $\frac{\mathrm{d}U}{\mathrm{d}\rho}$ im gesamten Intervall positiv ist.

Im zweiten Intervall ist $U$ stets positiv, $T$ zu Beginn positiv und am Ende negativ, dazwischen einmal $T=0$, und zwar nur einmal, weil $\frac{\mathrm{d}T}{\mathrm{d}\rho}$ im gesamten Intervall negativ ist.

Im dritten Intervall ist $T$ stets negativ, während $U$ einen Vorzeichenwechsel durchläuft, sodass \emph{einmal} $U=0$ gilt.

Im vierten Intervall ist $U$ stets negativ und $T$ \emph{einmal} $=0$.

In den folgenden Intervallen wiederholen sich diese Verhältnisse in derselben Reihenfolge, sodass das fünfte dem ersten, das sechste dem zweiten usw. entspricht.

